% Page 8: virtual state and authoritative ownership.
\subsection{虚拟状态结构与权威归属}
虚拟体验围绕有限的内置人物展开，纸条、聊天、待回复和周期回望常需一起呈现，
因此按账号组织一份互动状态，便于用户再次进入时继续同一段体验。
个人资料由账号统一管理，真人通信则独立保存。
表\ref{tab:design-state}按业务用途说明各组状态在互动流程中的作用。

\begin{table}[H]
\centering\small\linespread{1.15}\selectfont
\caption{虚拟互动状态的业务分组与作用}\label{tab:design-state}
\begin{tabularx}{\linewidth}{@{}p{.27\linewidth}p{.17\linewidth}X@{}}
\toprule 业务分组 & 处理阶段 & 作用与用户结果 \\ \midrule
体验进度 & 开始与继续 & 记录体验起点和当前进度，使纸条、交流与回望沿同一过程推进 \\
个人资料 & 进入页面 & 读取账号中的介绍、话题与开关，保持各场景对本人资料的理解一致 \\
纸条记录 & 发起与处理 & 保存与各人物的纸条及决定，区分等待处理、保留、接受和拒绝 \\
对话记录 & 持续交流 & 按人物保留历史消息，再次进入时接着阅读和交流 \\
待回复 & 发送后等待 & 区分聊天或纸条的回复任务，让页面说明尚未完成的互动 \\
屏蔽选择 & 停止互动 & 清除对应待回复与授权，保留已有纸条和对话供本人回看 \\
私人心笺 & 记录感受 & 保存本人对人物的感受与备注，供日常记录和回顾 \\
本人回望 & 表达选择 & 分别记录当期意愿与授权，使感受回顾和允许揭晓成为不同决定 \\
对方情境 & 虚拟演示 & 保存内置人物的当期选择，推动虚拟人物的周期互动 \\
揭晓记录 & 条件满足后 & 保存当期揭晓结果，每账号每周期一次机会，供后续回顾 \\
\bottomrule
\end{tabularx}
\end{table}

进入页面时组合账号资料与本人互动记录。首次使用建立初始体验，
再次进入继续读取历史，并补充新增人物的互动位置。
历史读取异常时提示重试，并保留原有互动记录。

用户提交操作后，服务端重新读取当前状态，检查该动作是否适合当前人物与阶段，再保存处理结果。
例如发出消息后先展示待回复，回复完成后进入对话历史；屏蔽人物时停止对应等待与授权，
已有纸条和聊天仍可保留。这样把“暂未完成”“停止继续”和“已有记录”分开表达。

资料归账号，虚拟互动归本人状态，相遇、媒体和真人消息各自管理，减少资料版本分歧。
关联保存与中断恢复见第\ref{sec:design-consistency}节。
心笺保存私人感受，回望记录当期选择；虚拟体验按演示进度推进，真人交流按实际日期计算额度。

\clearpage
